% chapters/theory/04.tex

\chapter{Training Large Language Models}
\label{chap:training-llms}

Large language models feel effortless at inference time—type a prompt, receive fluent text—but that apparent simplicity hides one of the most demanding training pipelines in modern machine learning. Training is expensive because it operates at extreme scale: billions of parameters, trillions of tokens, and massive GPU clusters. It is also technically intricate because memory, bandwidth, and numerical stability become first-class constraints. This chapter develops a practical understanding of the full training pipeline, from pretraining to supervised fine-tuning, and the engineering techniques that make it feasible.

\section{From task-specific training to pretrain-and-tune}
For many years, the dominant recipe in applied machine learning was task-specific: collect labeled data for a single objective, train a model from scratch, validate, and deploy. This is still valid, but it is inefficient for language. Most language tasks share an underlying skill—understanding and producing text—so it is wasteful to relearn language structure for every new application.

The modern recipe is transfer learning. First, we train a general-purpose model on vast amounts of unlabeled text so it learns the patterns of language and code. Second, we refine it for a narrower objective, such as behaving like a helpful assistant or solving a domain-specific task. In the LLM context these two phases are commonly called \emph{pretraining} and \emph{fine-tuning}. The key idea is that the expensive part—learning language—should be done once and reused.

\section{Pretraining: next-token prediction at internet scale}
Pretraining teaches a model the statistical structure of text by optimizing a simple objective: predict the next token given all previous tokens. If the token sequence is \(x_{1:t}\), the model learns to approximate \(P(x_{t+1}\mid x_{1:t})\). In decoder-only Transformers, causal masking allows the entire sequence to be processed in parallel during training while preserving the constraint that position \(t\) cannot use information from positions \(>t\).

The data used for pretraining is typically a mixture of sources that collectively approximate “written knowledge”: large web crawls (often labeled under the umbrella of Common Crawl), encyclopedic sources such as Wikipedia, discussion forums such as Reddit, and code corpora from GitHub and Stack Overflow-style sites. Because the training objective is general and self-supervised, pretraining can consume extremely large quantities of raw text and code without manual labeling. Scale is measured in tokens, and state-of-the-art training runs commonly reach hundreds of billions to trillions of tokens.

Pretraining is not merely expensive; it is also slow and operationally heavy. This has two important consequences. First, the trained model has a knowledge cutoff: it can only “know” what was present in the training data up to the point where the dataset was collected and frozen. Second, updating knowledge by continuing to train is difficult. Even if additional training injects new facts, it can cause regressions elsewhere—changing one region of the model’s behavior can unintentionally degrade another. For this reason, many real systems prefer to keep the base model fixed and use retrieval and tools at inference time for freshness rather than continuously modifying weights.

\section{Compute language: FLOPs versus FLOPS}
Discussions about LLM training often use two similar-looking terms that mean different things. FLOPs (floating-point operations) is a measure of how much computation is performed, while FLOPS (floating-point operations per second) is a measure of how fast hardware can perform those operations.

FLOPs helps reason about total training cost: as a rough mental model, the training compute grows with the number of tokens processed and the number of parameters that are active during computation. The exact constant depends on architecture, but the key relationship is intuitive: more parameters and more tokens imply more arithmetic. FLOPS, on the other hand, appears in hardware specifications: GPUs advertise throughput for FP32, FP16, BF16, and so on. Throughput matters because it determines wall-clock time given a fixed FLOP budget.

\section{Scaling laws and compute-optimal training}
Empirically, language model performance improves as we scale model size, dataset size, and compute. These observations led to a wave of scaling in which teams simply built larger models, often finding consistent performance gains. A more subtle question follows: if compute is fixed and expensive, how should we allocate it between model parameters and training tokens?

A key takeaway from compute-optimal analyses is that many early large models were “undertrained” relative to their parameter count: they were too large for the amount of data they consumed. A widely used rule-of-thumb that emerged is that a compute-efficient regime often trains on a number of tokens that is on the order of twenty times the number of parameters. The precise coefficient is not universal, but the core idea is: if you increase parameter count without increasing data accordingly, you do not fully realize the potential of the capacity you built.

In practice, major training efforts do not rely solely on a global rule. Teams typically run smaller-scale experiments—smaller models on smaller datasets—to estimate their own compute-optimal frontier, then extrapolate. This is especially relevant when architectural choices differ (for instance, mixture-of-experts models change how many parameters are active per token).

\section{The training loop and why memory dominates}
Training a Transformer is conceptually simple: forward pass, compute loss, backward pass, update weights. The difficulty is that each step requires storing large intermediate tensors.

During the forward pass, the model produces activations at every layer. These activations are required later to compute gradients during backpropagation. The loss is typically cross-entropy for next-token prediction. During the backward pass, gradients for parameters must be computed and stored. Finally, the optimizer update step needs additional state. Adam, for example, maintains moving averages of gradients and squared gradients (often called first and second moments). These optimizer states can be comparable in size to the parameters themselves.

In other words, training must store not only the parameters, but also gradients, optimizer states, and activations. Even the best GPUs have limited memory (tens of gigabytes). As a result, LLM training is often memory-bound rather than compute-bound, and many training optimizations exist primarily to fit the model and its training state into available memory.

\section{Distributing training: data parallelism and ZeRO}
A natural response to limited memory is to use many GPUs. The simplest strategy is data parallelism: split the batch across devices, replicate the model on each device, compute gradients locally, and then aggregate gradients across devices so that all replicas apply the same update. Data parallelism reduces the per-device batch size, which can reduce activation memory, but it introduces communication overhead due to gradient synchronization.

ZeRO (Zero Redundancy Optimizer) addresses a different inefficiency: in vanilla data parallelism, every GPU stores a complete copy of parameters, gradients, and optimizer states. ZeRO reduces duplication by sharding these quantities across GPUs. In its simplest form it shards optimizer state; in more advanced forms it shards gradients and then parameters as well. The result is a dramatic reduction in memory per GPU, which can be the difference between “fits” and “does not fit.” The trade-off is increased communication and orchestration complexity, because GPUs must exchange shards during forward and backward passes.

\section{Distributing training: model parallelism}
Data parallelism assumes each GPU can store a full model replica. When the model is too large, we must split the model itself across devices. This is model parallelism, and it comes in several common flavors.

Tensor parallelism splits large matrix multiplications across GPUs. Pipeline parallelism assigns different layers to different GPUs and streams micro-batches through the pipeline. Expert parallelism applies to mixture-of-experts architectures by placing different experts on different GPUs and routing tokens to the correct devices. In modern training systems these strategies are not mutually exclusive; large-scale training stacks multiple forms of parallelism together, combining data parallelism with tensor/pipeline/expert parallelism to meet both memory and throughput constraints.

\section{FlashAttention: exact attention, less memory traffic}
Attention is expensive not only because of arithmetic but because naïve implementations move large intermediate matrices in and out of GPU memory. GPUs have a memory hierarchy: a large but relatively slower memory (HBM) and a much smaller but faster on-chip memory (SRAM). The bottleneck in practice is often the repeated reads and writes to HBM required by the attention computation, especially around the softmax step.

FlashAttention is a family of methods that makes attention faster without approximating it. The key idea is tiling: rather than computing the entire \(QK^\top\) matrix and storing it, the algorithm processes attention in blocks that fit in SRAM. It streams blocks of \(Q\), \(K\), and \(V\) through SRAM, computes the necessary softmax-normalized partial results, and writes out only the final outputs. This dramatically reduces HBM traffic, which is often the true limiting factor.

A second important idea is recomputation. Backpropagation needs activations; storing them is expensive. If attention becomes fast enough, it can be better to store fewer activations and recompute them during the backward pass. Remarkably, this can reduce memory usage and improve runtime simultaneously, because the cost saved in HBM traffic outweighs the extra arithmetic. Later FlashAttention versions refine these ideas to match newer GPU architectures.

\section{Quantization and mixed precision}
Training systems also exploit the fact that not every computation requires full FP32 precision. Floating-point formats differ in how they represent exponent and mantissa bits, trading range and precision against memory and throughput. Lower precision formats such as FP16 and BF16 use half the memory of FP32 and often run faster on modern hardware.

Quantization generalizes this idea by mapping values to lower-precision representations, potentially reducing memory and compute at the cost of precision. Different quantization schemes exist (e.g., scale + zero-point versus absmax); not all layers are equally sensitive. The central question is always the same: how much numerical precision is actually needed to preserve training stability and downstream performance?

Mixed precision training is the workhorse strategy. A common pattern is to keep a master copy of weights in FP32 while running forward and backward computations in FP16 or BF16. Updates are applied to the FP32 master weights, reducing error accumulation. The intuition is that activations and gradients can tolerate lower precision noise, while weight updates benefit from higher precision to prevent drift and instability.

\section{From pretraining to supervised fine-tuning}
Pretraining produces a strong continuation model: it can autocomplete text, but it is not necessarily a helpful assistant. If you ask a question, a purely pretrained model may produce a plausible continuation rather than a useful response. Supervised fine-tuning (SFT) addresses this by training on curated input-output pairs that reflect the desired behavior.

In SFT, the prompt is treated as conditioning context rather than part of the prediction target. Loss is computed on the assistant response tokens, not on the user instruction tokens. Instruction tuning is a prominent form of SFT designed to convert the model into a broadly helpful assistant. The data mixture commonly includes tasks such as explanation, summarization, planning, story writing, code assistance, and structured formatting. Safety examples are also part of many instruction datasets, teaching the model to refuse harmful requests or respond with appropriate caution.

SFT datasets are far smaller than pretraining corpora, but they are much higher quality. This is why the tuning stage can be tractable even when pretraining is not.

\section{LoRA: parameter-efficient fine-tuning}
Full fine-tuning updates every weight, which is expensive in memory and compute. LoRA (Low-Rank Adaptation) reduces this cost by freezing a base weight matrix \(W_0\) and learning a low-rank update:
\[
W = W_0 + BA,
\]
where the rank \(r\) is small compared to the full dimension. Only the LoRA parameters \(A\) and \(B\) are trained. This can reduce trainable parameters by orders of magnitude, making fine-tuning feasible on smaller hardware and enabling rapid adaptation for multiple downstream tasks.

Early LoRA recipes often focused on attention projections, but later practice found that applying LoRA in feedforward blocks can be highly effective as well. In modern usage, LoRA is frequently placed in both attention and feedforward components, with the exact placement tuned empirically for the task.

\section{QLoRA: quantized base weights plus LoRA adapters}
QLoRA combines two ideas: keep the base model frozen but quantized to reduce VRAM, and train LoRA adapters in higher precision. The base weights are stored in a compact representation (notably NF4 in the original proposal, motivated by the approximate normality of weight distributions), while the LoRA updates are learned in BF16 or similar formats. A further optimization called double quantization compresses the quantization constants themselves. The practical impact is that models that would otherwise be too large to fine-tune on a given GPU can become feasible.

\section{Summary}
Training LLMs is best understood as a staged pipeline constrained by compute and memory. Pretraining is the dominant cost and teaches broad language and code structure through next-token prediction on massive corpora. Compute-optimal scaling emphasizes that model size must be matched with sufficient training tokens. The training loop is memory-heavy because it must store parameters, optimizer state, gradients, and activations. Practical training therefore relies on distributed strategies such as data parallelism, ZeRO sharding, and model parallelism, alongside hardware-aware kernels such as FlashAttention. After pretraining, supervised fine-tuning shapes behavior for instruction-following. Finally, parameter-efficient tuning methods such as LoRA and QLoRA make adaptation practical by dramatically reducing the number of trainable parameters and the memory footprint of fine-tuning.

\section{Exercises}
\paragraph{1. FLOPs vs FLOPS.}
Explain the difference between FLOPs and FLOPS in one paragraph. Give one example of how each is used in practice when discussing LLM training.

\paragraph{2. Why attention becomes memory-bound.}
Describe why naïve attention implementations spend significant time moving data rather than doing arithmetic. How does tiling reduce this bottleneck?

\paragraph{3. ZeRO intuition.}
Why is vanilla data parallelism redundant? Explain what quantities ZeRO shards and why that reduces memory. What cost does this introduce?

\paragraph{4. Mixed precision motivation.}
Why is it useful to keep master weights in FP32 while computing activations/gradients in FP16/BF16? What failure mode does this help prevent?

\paragraph{5. LoRA scaling.}
Assume a weight matrix \(W_0 \in \mathbb{R}^{d\times k}\). If LoRA uses rank \(r\), how many trainable parameters are introduced by \(A\) and \(B\)? Compare this to full fine-tuning.
